% Title: Tensor Product Representations and Analytic Continuation on Riemannian Structures
% Author: Amlal El Mahrouss (A. E. Mahrouss)

\documentclass[10pt]{article}

%% ---- Core Packages ----------------------------------------
\usepackage[T1]{fontenc}
\usepackage[utf8]{inputenc}
\usepackage{lmodern}
\usepackage{microtype}

%% ---- Page Geometry ----------------------------------------
\usepackage[
  top=1in, bottom=1in,
  left=0.75in, right=0.75in,
  columnsep=0.25in
]{geometry}

%% ---- Mathematics ------------------------------------------
\usepackage{amsmath, amssymb, amsthm, amsfonts}
\usepackage{mathtools}
\usepackage{bm}
\usepackage{textcomp}
\usepackage{esint}

%% ---- Figures & Tables -------------------------------------
\usepackage{graphicx}
\usepackage{booktabs}
\usepackage{multirow}
\usepackage{array}
\usepackage{caption}
\usepackage{subcaption}
\usepackage{float}

%% ---- Code Listings ----------------------------------------
\usepackage{listings}
\usepackage{xcolor}
\usepackage{algorithm}
\usepackage{algpseudocode}

\lstdefinestyle{codestyle}{
  backgroundcolor=\color{gray!8},
  basicstyle=\ttfamily\footnotesize,
  breakatwhitespace=false,
  breaklines=true,
  captionpos=b,
  commentstyle=\color{green!50!black},
  keywordstyle=\color{blue}\bfseries,
  stringstyle=\color{orange!70!black},
  numberstyle=\tiny\color{gray},
  numbers=left,
  numbersep=5pt,
  showstringspaces=false,
  frame=single,
  rulecolor=\color{gray!40},
  tabsize=2
}
\lstset{style=codestyle}

%% ---- Hyperlinks -------------------------------------------
\usepackage[
  colorlinks=true,
  linkcolor=blue!70!black,
  citecolor=green!50!black,
  urlcolor=blue!60!black
]{hyperref}
\usepackage{url}

%% ---- Bibliography -----------------------------------------
\usepackage[numbers, sort&compress]{natbib}

%% ---- Miscellaneous ----------------------------------------
\usepackage{enumitem}
\usepackage{siunitx}
\usepackage{cleveref}

%% ---- Theorem Environments ---------------------------------
\theoremstyle{plain}
\newtheorem{theorem}{Theorem}[section]
\newtheorem{lemma}[theorem]{Lemma}
\newtheorem{corollary}[theorem]{Corollary}
\newtheorem{proposition}[theorem]{Proposition}

\theoremstyle{definition}
\newtheorem{definition}[theorem]{Definition}
\newtheorem{example}[theorem]{Example}

\theoremstyle{remark}
\newtheorem{remark}[theorem]{Remark}

%% ---- Custom Commands --------------------------------------
\newcommand{\R}{\mathbb{R}}
\newcommand{\N}{\mathbb{N}}
\newcommand{\Z}{\mathbb{Z}}
\newcommand{\E}{\mathbb{E}}
\newcommand{\Prob}{\mathbb{P}}
\newcommand{\norm}[1]{\left\lVert #1 \right\rVert}
\newcommand{\abs}[1]{\left\lvert #1 \right\rvert}
\newcommand{\inner}[2]{\left\langle #1,\, #2 \right\rangle}
\newcommand{\ie}{\textit{i.e.}\xspace}
\newcommand{\eg}{\textit{e.g.}\xspace}
\newcommand{\etal}{\textit{et al.}\xspace}

%% ---- Caption Formatting -----------------------------------
\captionsetup{
  font=small,
  labelfont=bf,
  format=hang,
  justification=justified
}

\usepackage{lettrine}

%% ---- Header / Footer --------------------------------------
\usepackage{fancyhdr}
\pagestyle{fancy}
\fancyhf{}
\fancyhead[R]{\small\thepage}
\renewcommand{\headrulewidth}{0.4pt}

%% ============================================================
%%  FRONT MATTER
%% ============================================================

\title{%
	\vspace{-1.5em}%
	\large\textbf{Special use cases of the Cauchy Momentum Equation.}%
}
\author{By Amlal El Mahrouss}
\date{\today}

%% ============================================================
\begin{document}
%% ============================================================

\maketitle
\tableofcontents
\newpage

%% ============================================================
\section{Abstract:}
%% ============================================================

\lettrine[lines=1, lhang=0.1, loversize=0.2]{I}N this paper we will explore several possiblilities of the Cauchy momentum equation as stated below:

%% ============================================================
\section{Introduction and Definitions}
%% ============================================================

\begin{definition}
There exist by setting the cauchy stress tensor $\sigma$ to be the sum of a viscosity term $\tau$ and a pressure term $-p\, \mathbf{I}$, defined as the volumetric stress, the following equation:
\begin{equation}
\label{eq:cauchy}
  \rho \frac{\mathbf{D}\textbf{u}}{\mathbf{D}t} = - \nabla p + \nabla \cdot \tau + \rho \, \textbf{a}.
\end{equation}
\end{definition}
\begin{remark}
\label{rem:1}
  The material derivative is defined as: $\frac{\mathbf{D}\textbf{u}}{\mathbf{D}t}$.
\end{remark}

\begin{remark}
\label{rem:2}
  The density variable is defined as: $\rho$.
\end{remark}

\begin{remark}
\label{rem:3}
  The flow density variable is defined as: $\textbf{u}$.
\end{remark}

\begin{remark}
\label{rem:4}
  The divergence vector operator, defined as: $\nabla \cdot$.
\end{remark}

\begin{remark}
\label{rem:5}
  The pressure variable, defined as: $p$.
\end{remark}

\begin{remark}
\label{rem:6}
  The time variable, defined as: $t$.
\end{remark}

\begin{remark}
\label{rem:7}
  The deviatoric stress tensor operator, defined as: $\tau$.
\end{remark}
\begin{remark}
  The body accelerations defined as $\mathbf{a}$ acting on the continuum.
\end{remark}

\subsection{Several Applications from Tensors.}

Assuming we use the defintion per Remark~\ref{rem:7}, Redefinining Equation~\ref{eq:cauchy}, that yields:
\begin{equation}
  \rho \frac{\mathbf{D}\textbf{u}}{\mathbf{D}t} = - \nabla p + \nabla \cdot \tau + \rho \, \textbf{a}.
\end{equation}
Assuming the variable of integration to be Remark~\ref{rem:7} between the open intveral $\mathbb{R}^{+} \land +\infty$:
\begin{equation}
  \int_{\mathbb{R}^{+}}^{+\infty} \rho \frac{\mathbf{D}\textbf{u}}{\mathbf{D}t} \, d\tau = - \int_{\mathbb{R}^{+}}^{+\infty} \nabla p + \nabla \cdot \tau + \rho \, \textbf{a} \, d\tau.
\end{equation}

%% ============================================================
\section{Conclusion:}
%% ============================================================

%% TODO: Rework the references at the end.
\nocite{*}
\bibliographystyle{acm}
%\bibliography{}

%% ============================================================
\end{document}
%% ============================================================